% !PRINTED 516
\weberpage{516}

Diese Resultate sind von grosser Wichtigkeit und verdienen
besonders hervorgehoben zu werden.

Nach \gesperrt{Kronecker} heissen, wenn $F(x) = 0$ die gegebene
Gleichung ist und $N$ der zugehörige \gesperrt{Galois}'sche Körper, die
Grössen von $N$ \gesperrt{die der Gleichung $F(x) = 0$ natürlichen
Irrationalitäten}.

Wir haben dann den Satz:

\gesperrt{Jede mögliche Reduction der Galois'schen Gruppe wird
herbeigeführt durch Adjunction einer natürlichen
Irrationalität. Ist $j$ der Index der reducirten Gruppe in
Bezug auf die ursprüngliche, so kann die Reduction nicht
eintreten durch Adjunction eines Körpers von niedrigerem als
dem $j^{\mathrm{ten}}$ Grade. Der Grad kann aber gleich $j$ sein
und dann ist die adjungirte Irrationalität eine natürliche. Ist
der Grad des adjungirten Körpers grösser als $j$, so ist er ein
Vielfaches von $j$.}

\begin{center}
\S.\,158.

\gesperrt{Imprimitive Gruppen.}
\end{center}

Wir machen von unseren allgemeinen Sätzen die Anwendung auf
Gruppen von besonderer Natur. Zunächst nehmen wir die
\gesperrt{imprimitiven Gruppen}. Wir haben im \S.\,151 gesehen, dass, wenn
$f(x) = 0$ eine irreducibele Gleichung $n^{\mathrm{ten}}$
Grades mit imprimitiver Gruppe ist, sich $n$ so in zwei
Factoren $r \boldsymbol{.} s$ zerlegen lässt, dass die Wurzeln in
$s$ Reihen von je $r$ Gliedern zerfallen:

\weq{1}{\begin{aligned}
A &= \alpha,\ \alpha_1 \dots \alpha_{r-1},\\
B &= \beta,\ \beta_1 \dots \beta_{r-1},\\
&\cdots\cdots\cdots\cdots\cdots\cdots\cdots\cdots\\
S &= \sigma,\ \sigma_1 \dots \sigma_{r-1},
\end{aligned}}

und dass durch die Permutationen der Gruppe $P$ die Grössen der
einzelnen Reihen nicht von einander getrennt, sondern nur unter
einander permutirt werden, während andererseits die Reihen
wieder unter sich permutirt werden.

Unter den Permutationen der \gesperrt{Galois}'schen Gruppe $P$ von $f(x)$
wird es nun gewisse geben, deren Inbegriff wir mit $Q$
%CONT>>
